\section{Introduction}
A document assistant must keep answering questions as its sources change.
New documents need to become searchable, and answers need to be supported by
the retrieved text. Both requirements have been studied in search and RAG
systems \citep{freshdiskann,metis,alce}. When ingestion and LLM inference run
on the same small machine, they also compete for CPU time and memory.

There are several ways to share those resources. Delaying an update gives
inference more time but makes users wait for new information. Reserving cores
for ingestion can help it finish sooner. A fixed reservation, however, can
continue to limit inference after embedding has finished. We study how these
choices affect the whole request. Does an earlier index update lead to an
earlier answer, or does the cost appear elsewhere?

Recent measurement studies ask similar questions about other parts of LLM
serving. Breaking the Ice identifies the stages responsible for startup time.
\emph{Speculative Decoding: Performance or Illusion?} checks whether speedups
seen in prototypes hold in a production engine \citep{breakingice,specillusion}.
\emph{Beyond the Buzz} studies how workloads and configurations change the
benefits of disaggregation \citep{beyondbuzz}. These works show the value of
measuring where an optimization saves time and where it adds cost.

Scheduling systems already account for different inference phases.
Sarathi-Serve uses chunked prefills, DistServe separates prefill and decode
resources, and SuperInfer coordinates scheduling with KV transfers
\citep{sarathi,distserve,superinfer}. METIS adapts RAG configuration and
scheduling to meet quality requirements \citep{metis}. We build on this
perspective with an evaluation framework for CPU allocation, document publication and
answer completion on a Raspberry Pi 5. We evaluate fixed allocations.
We do not introduce a new adaptive scheduler.

Our method has three steps. We first run inference and ingestion separately
and together, recording phase times and checking which CPUs each thread can
use. We then run matched application episodes with different update policies
and request arrival times. Finally, we fix timing predictions before testing
them on other documents and arrivals. At every step, we check answer quality,
citations and the retrieved version separately from response time.
Section~\ref{sec:problem} defines the evaluation steps.

We run the full Qwen2.5-14B-Instruct model in Q4\_K\_M on a Raspberry Pi 5
with 16 GB RAM, alongside the MiniLM document worker. This gives a concrete
setting for studying contention on a small CPU system. The 27 diagnostic
cells show a clear trade-off. Reserving two
cores for ingestion almost halves update time under concurrency, but increases
the long-prompt answer time from 174.3 to 246.4\,s. Most of the extra time is in
prefill. In 36 application episodes, publication finishes before current
inference under both concurrent policies, even when the index is rebuilt.
The fixed partition also slows the next request. In 12 further episodes, the
predictions correctly rank the two policies in six pairs, but their accuracy
is no better than a simple baseline added after the experiment.

Our contribution is an evaluation methodology and its implemented framework.
We built the controller for CPU allocation and independent request arrivals,
connected versioned document publication to retrieval and inference, and added
phase measurements, resource sampling and answer checks. The framework retains
inputs, outputs and prediction hashes so each comparison can be traced to its
configuration. These components make the allocation trade-offs measurable with
existing runtimes. Earlier precision, backend and QA experiments provide
separate controls rather than extra repetitions of the 75 cells. Our documents
and arrivals are controlled cases, not production traces. Dynamic core
reassignment and independent human review of answer correctness remain open.
